% ============================================================================
%  Informe corto para el docente - Taller Práctico II
%  Compilar:  latexmk -pdf informe_docente.tex
%  Usa los mismos estilos (config.tex) y el mismo firmware que el informe largo.
% ============================================================================
\documentclass[11pt,letterpaper]{article}

\usepackage[utf8]{inputenc}
\usepackage[T1]{fontenc}
\usepackage[spanish,es-tabla,es-nodecimaldot]{babel}
\usepackage{lmodern}
\usepackage{microtype}
\usepackage[margin=2.2cm]{geometry}
\usepackage{fancyhdr}
\usepackage{titlesec}
\usepackage{enumitem}
\setlist{itemsep=1pt, topsep=3pt}
\usepackage{amsmath,amssymb}
\usepackage{siunitx}
\sisetup{output-decimal-marker={.}, group-separator={\,}, per-mode=symbol}
\usepackage{booktabs}
\usepackage{tabularx}
\usepackage{array}
\usepackage{float}
\usepackage{graphicx}
\usepackage{caption}
\captionsetup{font=small, labelfont=bf, skip=4pt}
\usepackage{tikz}
\usetikzlibrary{arrows.meta, positioning, calc, fit, babel}
\usepackage{xcolor}
\usepackage{listings}
\usepackage[hidelinks]{hyperref}

% ============================================================================
%  config.tex - Datos del informe, colores, estilos y macros
% ============================================================================

% ---------- Datos del documento (editar aquí) ----------
\newcommand{\titulo}{Robot móvil diferencial con PIC16F877A y DRV8833}
\newcommand{\subtitulo}{Taller Práctico II: Control dual de motores y robótica móvil diferencial}
\newcommand{\asignatura}{Microcontroladores y Sistemas Embebidos}
\newcommand{\programa}{Programa de Ingeniería Electrónica}
\newcommand{\facultad}{Facultad de Ingeniería}
\newcommand{\universidad}{Universidad de Ibagué}
\newcommand{\docente}{\pendiente{Nombre del docente}}
\newcommand{\fecha}{Noviembre de 2026}
\newcommand{\versiondoc}{0.2 (borrador)}

% ---------- Colores ----------
\definecolor{acento}{RGB}{0,0,0}          % titulos y reglas en negro
\definecolor{acentoclaro}{RGB}{236,236,236} % relleno gris claro en figuras
\definecolor{gris}{RGB}{110,110,110}
\definecolor{codigofondo}{RGB}{247,247,247}
\definecolor{codigocomentario}{RGB}{100,100,100}
\definecolor{codigoclave}{RGB}{0,0,0}
\definecolor{codigocadena}{RGB}{60,60,60}
\definecolor{pendientecolor}{RGB}{200,0,0}

% ---------- Marcas de contenido pendiente ----------
% Uso: \pendiente{qué falta}. Buscar "\pendiente" para ver lo que falta.
\newcommand{\pendiente}[1]{\textcolor{pendientecolor}{\textit{[#1]}}}

% Recuadro para fotografías o capturas que aún no se han agregado.
% Uso: \figurapendiente{descripción}{etiqueta}
\newcommand{\figurapendiente}[2]{%
  \begin{figure}[H]
    \centering
    \fbox{\parbox[c][4cm][c]{0.8\textwidth}{\centering\pendiente{Agregar: #1}}}
    \caption{#1}
    \label{#2}
  \end{figure}}

% ---------- Encabezado y pie ----------
\pagestyle{fancy}
\fancyhf{}
\fancyhead[L]{\small\color{gris}\asignatura}
\fancyhead[R]{\small\color{gris}\universidad}
\fancyfoot[L]{\small\color{gris}Versión \versiondoc}
\fancyfoot[R]{\small\thepage}
\renewcommand{\headrulewidth}{0.4pt}
\setlength{\headheight}{14pt}

% ---------- Títulos de sección ----------
\titleformat{\section}{\Large\bfseries\color{acento}}{\thesection.}{0.6em}{}
\titleformat{\subsection}{\large\bfseries\color{acento}}{\thesubsection}{0.6em}{}
\titlespacing*{\section}{0pt}{18pt}{8pt}
\titlespacing*{\subsection}{0pt}{12pt}{6pt}

% ---------- Tablas ----------
\renewcommand{\arraystretch}{1.25}
\newcolumntype{L}{>{\raggedright\arraybackslash}X}
\newcolumntype{C}{>{\centering\arraybackslash}X}

% ---------- Listados de código ----------
\lstdefinestyle{c-xc8}{
  language=C,
  basicstyle=\ttfamily\footnotesize,
  backgroundcolor=\color{codigofondo},
  commentstyle=\color{codigocomentario}\itshape,
  keywordstyle=\color{codigoclave}\bfseries,
  stringstyle=\color{codigocadena},
  morekeywords={__interrupt, volatile, __delay_ms},
  numbers=left,
  numberstyle=\tiny\color{gris},
  numbersep=8pt,
  frame=single,
  rulecolor=\color{gris!40},
  breaklines=true,
  breakatwhitespace=true,
  showstringspaces=false,
  tabsize=4,
  xleftmargin=1.5em,
  framexleftmargin=1.5em,
  captionpos=t
}
\lstset{style=c-xc8}
\renewcommand{\lstlistingname}{Código}
\renewcommand{\lstlistlistingname}{Índice de códigos}

% ---------- Metadatos del PDF ----------
\hypersetup{
  pdftitle={\titulo},
  pdfsubject={\subtitulo},
  pdfkeywords={PIC16F877A, DRV8833, PWM, robótica móvil, XC8}
}

\titlespacing*{\section}{0pt}{12pt}{5pt}
\titlespacing*{\subsection}{0pt}{8pt}{4pt}
\renewcommand{\arraystretch}{1.15}

\begin{document}

% ---------------------------- Encabezado ------------------------------------
\begin{center}
  {\small\scshape \universidad\ --- \programa\par}
  \vspace{2pt}
  {\small \asignatura\par}
  \vspace{8pt}
  {\LARGE\bfseries\color{acento} Taller Práctico II: control dual de motores\par}
  \vspace{2pt}
  {\large Robot móvil diferencial con PIC16F877A y DRV8833\par}
  \vspace{8pt}
  {\small
  \textbf{Integrantes:} Brahiam Alberto Campiño Romero,
  \pendiente{integrante 2}, \pendiente{integrante 3}
  \quad\textbf{Docente:} \docente
  \quad\textbf{Fecha:} \fecha\par}
  \vspace{2pt}
  {\color{acento}\rule{\textwidth}{0.8pt}}
\end{center}

\noindent\textbf{Resultado.} Se configuró el PWM de hardware (CCP1 y CCP2 con
Timer2) a \SI{1250}{\hertz} con resolución de 10~bits y se implementó el
control de los dos motores con el DRV8833: avance, reversa, giro sobre su eje
y freno dinámico, con rampas de aceleración. Como extensión, el carro se maneja
desde un celular Android con un control tipo gamepad por Bluetooth (HC-05). En
la tarjeta del curso se encontró que el pin de habilitación del driver (EEP)
no está conectado, lo que impedía mover los motores con el código de la guía;
la corrección se describe en la sección~\ref{c:montaje}.

% ============================================================================
\section{Montaje físico y DRV8833}
\label{c:montaje}

\begin{table}[H]
  \centering
  \caption{Conexiones en la tarjeta (PIC16F877A en TQFP-44).}
  \label{c:tab:pines}
  \small
  \begin{tabularx}{\textwidth}{@{}l c l L@{}}
    \toprule
    \textbf{Pin PIC} & \textbf{Pata} & \textbf{Destino} & \textbf{Función} \\
    \midrule
    RC2 / CCP1 & 36 & MIPWM (IN2) & PWM del motor izquierdo \\
    RD0        & 38 & MIPIN (IN1) & Dirección del motor izquierdo \\
    RC1 / CCP2 & 35 & MDPWM (IN3) & PWM del motor derecho \\
    RD1        & 39 & MDPIN (IN4) & Dirección del motor derecho \\
    RC7 / RX, RC6 / TX & 1, 44 & HC-05 & Comandos desde el celular \\
    RB5, RB4, RB3 & 15, 14, 11 & LED1--LED3 & Diagnóstico \\
    MCLR, RB6, RB7 & 18, 16, 17 & ICSP & Programación con PICkit~3 \\
    \bottomrule
  \end{tabularx}
\end{table}

\begin{itemize}
  \item \textbf{Alimentación y masa común:} lógica a \SI{5}{\volt} con LM7805;
        motores desde la batería (VM $\leq$ \SI{10}{\volt}) con GND común cerca
        del DRV8833. Una pila de \SI{9}{\volt} no soporta el arranque de los
        motores (el PIC y el HC-05 se reinician); se usan dos baterías 18650
        (\SI{7.4}{\volt}).
  \item \textbf{Desacoplo:} \SI{0.1}{\micro\farad} en VDD--VSS del PIC;
        \SI{100}{\micro\farad} en paralelo con \SI{0.1}{\micro\farad} en
        VM--GND del driver; \SI{0.1}{\micro\farad} en los terminales de cada
        motor; cristal de \SI{20}{\mega\hertz} con \SI{22}{\pico\farad}.
  \item \textbf{ICSP activo:} MCLR, RB6 y RB7 quedan libres para el PICkit~3.
        El conector debe ir con la flecha del pin~1 hacia la marca de la
        tarjeta: conectado al revés aplica unos \SI{13}{\volt} en RB3/PGM, y así
        se dañó el PIC de una de las tarjetas.
  \item \textbf{nSLEEP (EEP):} en la tarjeta este pin no está conectado. El
        DRV8833 lo lleva a GND internamente, así que el driver queda dormido:
        se midieron \SI{0.14}{\volt}. Con una resistencia de
        \SI{10}{\kilo\ohm} a \SI{5}{\volt} se midieron \SI{1.41}{\volt}, porque
        el módulo tiene su propio \emph{pull-down} y se forma un divisor; ese
        valor está por debajo del umbral de \SI{2.5}{\volt} de la hoja de datos.
        Se conectó EEP directo a \SI{5}{\volt} (VCC del HC-05). En protoboard el
        firmware genera el pulso de activación por RD2.
\end{itemize}

\figurapendiente{Fotografía de la tarjeta montada}{c:fig:montaje}

% ============================================================================
\section{Justificación matemática}

\noindent\textbf{1. Ciclo de instrucción.}
\begin{equation*}
  T_{osc} = \frac{1}{\SI{20}{\mega\hertz}} = \SI{50}{\nano\second}
  \qquad
  T_{cy} = 4\,T_{osc} = \SI{200}{\nano\second}
  \qquad
  F_{cy} = \SI{5}{\mega\hertz}
\end{equation*}

\noindent\textbf{2. Registro PR2 (prescaler 1:16).}
\begin{equation*}
  PR2 = \frac{F_{osc}}{4\cdot \text{Pre}\cdot F_{PWM}} - 1
      = \frac{\num{20000000}}{4\cdot 16\cdot 1200} - 1 = 259.42
\end{equation*}
PR2 es de 8~bits (máximo 255), así que \SI{1.2}{\kilo\hertz} no es alcanzable:
la frecuencia mínima con este cristal y prescaler es
$\num{20000000}/(64\cdot 256) = \SI{1220.7}{\hertz}$. Se usa $PR2 = 249$:
\begin{equation*}
  F_{PWM} = \frac{\num{20000000}}{4\cdot 16\cdot 250} = \SI{1250}{\hertz}
  \qquad
  T_{PWM} = \SI{800}{\micro\second}
\end{equation*}
con un error de \SI{4.2}{\percent} frente a los \SI{1.2}{\kilo\hertz} pedidos.
Se elige 249 y no 255 porque deja exactamente 1000 pasos por periodo, y el
ciclo de trabajo se escribe directamente en escala 0--1000.

\noindent\textbf{3. Resolución del ciclo de trabajo (10 bits).}
\begin{gather*}
  T_{on} = DC\cdot T_{osc}\cdot\text{Pre} = DC\cdot\SI{0.8}{\micro\second}
  \qquad
  \text{pasos} = 4\,(PR2+1) = 1000 \\
  \text{resolución} = \log_2 1000 = 9.97~\text{bits}
  \qquad
  \Delta D = \SI{0.1}{\percent}
\end{gather*}
El valor de 10~bits se reparte en CCPR$x$L (8~bits altos) y CCP$x$X:CCP$x$Y
(2~bits bajos). Por ejemplo, $DC = 500$ da CCPR$x$L = 125, X:Y = 00 y
$T_{on} = \SI{400}{\micro\second}$ (50\,\%).

% ============================================================================
\section{Estructura del código XC8}

Se implementó el código de la guía y se amplió en un archivo propio,
organizado por capas:

\begin{table}[H]
  \centering
  \caption{Funciones principales.}
  \small
  \begin{tabularx}{\textwidth}{@{}>{\ttfamily\raggedright\arraybackslash}p{5.0cm} L@{}}
    \toprule
    \textbf{\normalfont Función} & \textbf{Qué hace} \\
    \midrule
    PWM\_Init() & PR2 = 249, CCP1/CCP2 en PWM, Timer2 con prescaler 1:16. \\
    Set\_Duty\_Motor\_Izq/\allowbreak Der(duty) & Carga el valor de 10~bits en CCPR$x$L y CCP$x$X:CCP$x$Y. \\
    Motor\_Izquierdo/\allowbreak Derecho(speed, dir) & Avance: dirección = 0 y PWM directo. Reversa: dirección = 1 y PWM invertido (1000 $-$ speed). \\
    Freno\_General() & $x$IN1 = $x$IN2 = 1 en ambos canales: freno dinámico. \\
    Mover(), Rampa\_Subir/\allowbreak Bajar() & Avance, reversa y giros con rampas de 0 a 1000 en pasos de 50 cada \SI{100}{\milli\second}. \\
    Procesar\_Byte(), \allowbreak Comando\_Gamepad() & Recepción por interrupción de los comandos del celular. \\
    \bottomrule
  \end{tabularx}
\end{table}

\noindent\textbf{Lógica de dirección y frenado.} Con la dirección en 0, el PWM
alterna entre avance (10) y parada libre (00). Con la dirección en 1, alterna
entre reversa (01) y freno (11); como el motor solo empuja mientras el PWM está
en 0, el ciclo se invierte para que la velocidad siga siendo proporcional al
valor pedido. Antes de invertir el sentido de una rueda se frena
\SI{30}{\milli\second} para limitar el pico de corriente.

\noindent\textbf{Cambios respecto al código de la guía:}
\begin{itemize}
  \item Se agregaron los giros a la izquierda y a la derecha sobre su eje, que
        la rúbrica evalúa y la secuencia base no tenía.
  \item Se agregó el pulso de activación de nSLEEP (por RD2, para el montaje
        en protoboard).
  \item El texto copiado del PDF impedía compilar (los encabezados de página
        quedaron dentro del código); se limpió.
  \item Control remoto con la app Dabble: cruceta, joystick proporcional
        (mezcla $v_{izq} = v_y + v_x$, $v_{der} = v_y - v_x$) y botones de
        velocidad. El PIC detecta solo la velocidad en baudios del HC-05 y
        frena si se pierde la señal.
  \item Watchdog y reinicio por bajo voltaje activos, para que una caída de la
        batería reinicie el PIC en vez de bloquearlo.
\end{itemize}
Los dos programas compilan con MPLAB XC8 sin errores. La lógica del control
remoto se verificó además en el computador con 22 pruebas automáticas.

% ============================================================================
\section{Demostración en hardware}

\begin{table}[H]
  \centering
  \caption{Resultados en la tarjeta.}
  \small
  \begin{tabularx}{\textwidth}{@{}l L c@{}}
    \toprule
    \textbf{Prueba} & \textbf{Observación} & \textbf{Resultado} \\
    \midrule
    PWM en RC2 y RC1 & Periodo \SI{800}{\micro\second} esperado; medido: \pendiente{} & Pendiente \\
    Avance y reversa con rampas & \pendiente{} & Pendiente \\
    Giro izquierda y derecha & \pendiente{} & Pendiente \\
    Freno sin reinicios del PIC & \pendiente{} & Pendiente \\
    Control desde el celular & Comandos recibidos y decodificados (verificado con los LEDs); falta la prueba con baterías 18650. & Parcial \\
    \bottomrule
  \end{tabularx}
\end{table}

\noindent\textbf{Problemas encontrados y solución:}
\begin{enumerate}
  \item \emph{Motores sin movimiento con el código de la guía:} EEP no
        conectado en la tarjeta $\rightarrow$ EEP a \SI{5}{\volt}.
  \item \emph{Comandos del celular mal interpretados:} el HC-05 no trabajaba a
        9600~baudios $\rightarrow$ detección automática de la velocidad.
  \item \emph{Reinicios y desconexiones al arrancar el motor:} la pila de
        \SI{9}{\volt} no da la corriente $\rightarrow$ baterías 18650 y freno
        automático por pérdida de señal.
  \item \emph{Falla al grabar una tarjeta:} PIC dañado por el ICSP conectado
        al revés $\rightarrow$ se trabajó con otra tarjeta.
\end{enumerate}

% ============================================================================
\section{Instrucciones de uso}

\noindent\textbf{Compilar.} En MPLAB X (v6.20, la última compatible con el
PICkit~3) con XC8, el proyecto tiene tres configuraciones: \texttt{try1}
(código de la guía), \texttt{completo} (código propio) y \texttt{leds} (prueba
de montaje). Se elige en el desplegable de la barra y se compila con
\emph{Clean and Build}; el archivo queda en
\texttt{dist/<configuración>/production/}. En el código propio,
\texttt{MODO\_BLUETOOTH = 0} ejecuta la secuencia del taller y
\texttt{MODO\_BLUETOOTH = 1} activa el control por celular.

\noindent\textbf{Grabar.} En la aplicación \emph{PICkit~3 Programmer}: conectar
el ICSP con la flecha hacia el pin~1, activar \emph{Manual Device Select},
elegir \emph{Midrange/Standard} y PIC16F877A, \emph{Import Hex}, \emph{Write}.
Con la batería conectada, \emph{VDD On} debe estar desactivado. Desmarcar
\emph{/MCLR} para que el programa arranque.

\noindent\textbf{Simular.} En Proteus, el DRV8833 se reemplaza por un L293D
(IN1 = RC2, IN2 = RD0, IN3 = RC1, IN4 = RD1; EN1 y EN2 a \SI{5}{\volt}). El PIC
se configura a \SI{20}{\mega\hertz} y el osciloscopio en RC2 y RC1 debe mostrar
un periodo de \SI{800}{\micro\second}. Las RPM positivas indican avance y las
negativas, reversa.

\noindent\textbf{Manejar desde el celular.} Vincular el HC-05 en los ajustes de
Bluetooth de un celular Android (clave 1234 o 0000). En la app \emph{Dabble}
(STEMpedia), con la ubicación activa, conectar el HC-05 y abrir
\emph{Gamepad}. Presionar varios botones hasta que el LED1 cambie con cada
pulsación. La cruceta mueve el carro, el joystick lo maneja de forma
proporcional, $\triangle$/$\times$ suben y bajan la velocidad y $\square$
frena. El HC-05 no funciona con iPhone; para iPhone se necesita un módulo BLE
(HM-10).

% ============================================================================
\section{Conclusiones}

\begin{itemize}
  \item Con \SI{20}{\mega\hertz} y prescaler 1:16 la frecuencia de
        \SI{1.2}{\kilo\hertz} no es alcanzable; $PR2 = 249$ da
        \SI{1250}{\hertz} y una escala exacta de 1000 pasos de 0.1\,\%.
  \item Una sola línea de dirección y un PWM por motor bastan para controlar
        velocidad, sentido y freno dinámico con el DRV8833.
  \item El fallo principal fue de hardware (nSLEEP sin conexión). Se encontró
        midiendo y comparando con la hoja de datos, no cambiando el código.
\end{itemize}

\noindent El código fuente, el informe técnico completo y las instrucciones
están en \url{https://github.com/Brahiam18/robot-diferencial-pic16f877a}.

\end{document}
